\chapter{Reproducing the Experiments}\label{app:reproduce}

The commands below reproduce the study from a fresh clone, in the order it was run. They assume a Linux machine with
two 24~GB GPUs, Docker, and Hugging Face access to the two Llama checkpoints. Every step resumes from its files, so an
interrupted command can simply be run again. Long commands should run in their own \texttt{tmux} session.

\section*{1. Environment and data}
\begin{lstlisting}[language=bash]
git submodule update --init                     # third_party/Table-Specialist at af1921f
python3.12 -m venv .venv && .venv/bin/pip install -r configs/requirements.lock
scripts/download_data.sh                         # WikiSQL, WikiTQ, AutoTest, BIRD dev, GitTables (96 zips)
# build the sandbox images from the repository's Dockerfiles and sandbox/Dockerfile.r with the base digests
# recorded in configs/sandbox.env, then:
set -a; . configs/sandbox.env; . configs/runtime.env; set +a
.venv/bin/python -m ts.corpus.build              # 10,053 tables, results/corpus_manifest.json
.venv/bin/python -m ts.bench.wikisql && .venv/bin/python -m ts.bench.wikitq \
  && .venv/bin/python -m ts.bench.spider && .venv/bin/python -m ts.bench.bird
\end{lstlisting}

\section*{2. NL-to-SQL core study and extensions}
\begin{lstlisting}[language=bash]
scripts/run_study.sh table_specialist 0 8000     # base TS, seeds 1-3, rounds 1-3
scripts/run_study.sh no_validation 1 8001        # base NV
.venv/bin/python -m ts.ext2 --seed 1 --gpu 0 --port 8000           # bestof4 + DPO (seeds 1-3)
.venv/bin/python -m ts.repair_ts --seed 1 --gpu 0 --port 8000      # repair-TS round 1 (seeds 1-3)
.venv/bin/python -m ts.repair_loop --seed 1                        # repair-TS rounds 2-3 (seeds 1-3)
TS_PROFILE=base .venv/bin/python -m ts.data_control --seeds 1 2 3
.venv/bin/python -m ts.sc_run dev && .venv/bin/python -m ts.sc_run test --temperature 1.0
# Instruct track: the same modules with TS_PROFILE=instruct, e.g.
TS_PROFILE=instruct .venv/bin/python -m ts.loop --arm table_specialist --seed 1 --gpu-pool
\end{lstlisting}

\section*{3. Other tasks}
\begin{lstlisting}[language=bash]
TS_PROFILE=instruct .venv/bin/python -m ts.eval.scala_eval generate
TS_PROFILE=instruct .venv/bin/python -m ts.eval.scala_eval score
TS_PROFILE=instruct .venv/bin/python -m ts.eval.ed_eval            # vanilla error detection
TS_PROFILE=instruct .venv/bin/python -m ts.ed_loop --seed 1        # seeds 1-3
TS_PROFILE=instruct .venv/bin/python -m ts.repair_ed --variant repair --seeds 1 2 3
TS_PROFILE=instruct .venv/bin/python -m ts.repair_ed --variant resample --seeds 1 2 3 --rounds 1
TS_PROFILE=instruct .venv/bin/python -m ts.repair_ed --variant samedata --seeds 1 2 3
TS_PROFILE=instruct TS_R_GPUS=0 TS_R_HOLD_GPU=1 .venv/bin/python -m ts.r_loop --arm ts --seeds 1 2 3
TS_PROFILE=instruct TS_R_GPUS=1 TS_R_HOLD_GPU=1 .venv/bin/python -m ts.r_loop --arm repairts --seeds 1 2 3
TS_PROFILE=instruct .venv/bin/python -m ts.eval.r_eval generate     # all 19 models, then:
TS_PROFILE=instruct .venv/bin/python -m ts.eval.r_eval score
\end{lstlisting}

\section*{4. Analysis and this report}
\begin{lstlisting}[language=bash]
.venv/bin/python -m ts.analysis.significance && .venv/bin/python -m ts.analysis.significance --task scala
.venv/bin/python -m ts.analysis.significance --task r && .venv/bin/python -m ts.analysis.ed_significance
.venv/bin/python -m ts.analysis.results_table > results/analysis/results_table.md
report/build.sh            # regenerates every table, then compiles main.pdf in a pinned TeX Live image
\end{lstlisting}

\section*{What cannot be reproduced exactly}
\begin{itemize}
\item vLLM's batched greedy decoding is not guaranteed to be bit-identical across runs and hardware. Our generations
are cached in \fpath{runs/llm_cache.sqlite}, which makes reruns on this machine exact.
\item The paper's own corpus, case lists and settings are unpublished (\cref{ch:limitations}).
\end{itemize}
